\subsection{A hyperplane Green envelope on the ball}
\label{subsec:ball-divisor-green}\index{Green function!hypersurface pole}

For a logarithmic pole along a hyperplane section of the unit ball, the envelope can be computed explicitly by restricting to the complex lines normal to the hyperplane. The example also shows how envelopes of bounded obstacles recover the finite part of a logarithmic pole, and that continuity inside the domain does not imply continuity up to the boundary.

In this subsection let $n\geq2$ and $m>0$, and write
\[
 B=\{(\zeta,w)\in\C\times\C^{n-1}:|\zeta|^2+|w|^2<1\},
 \qquad D=B\cap\{\zeta=0\}.
\]
Define $G_m$ as the pointwise supremum of the functions $u\in\PSH(B)$ such that $u\leq0$ and, near every point of $D$, $u\leq m\log|\zeta|^2+C$ for some local constant $C$. In particular, $u=-\infty$ on $D$. The constants need not stay bounded as the point of $D$ approaches $\partial B$.

\Needspace{9\baselineskip}
\begin{proposition}\label{prop:ball-divisor-green}
The envelope $G_m$ is given by
\begin{equation}\label{eq:ball-divisor-green}
 G_m(\zeta,w)=m\log\frac{|\zeta|^2}{1-|w|^2}.
\end{equation}
It is plurisubharmonic and continuous on $B$ with values in $[-\infty,0)$. The finite part $G_m-m\log|\zeta|^2$, first defined off $D$, extends smoothly across $D$ as the plurisubharmonic function
\[
 H_m(w)=-m\log(1-|w|^2).
\]
\end{proposition}
\begin{proof}
The function $-\log(1-|w|^2)$ is smooth and plurisubharmonic on $|w|<1$: its complex Hessian is
\[
 \left(\frac{\delta_{jk}}{1-|w|^2}
       +\frac{\overline w_jw_k}{(1-|w|^2)^2}\right)_{j,k}.
\]
Thus the right-hand side of \eqref{eq:ball-divisor-green} is plurisubharmonic; it is negative on $B$ and satisfies the pole condition, so it is a competitor. For the reverse inequality, fix a competitor $u$ and a point $w$ with $|w|<1$. The restriction $u(\cdot,w)$ is either identically $-\infty$ or subharmonic on the disk of radius $R=(1-|w|^2)^{1/2}$. In the latter case, $u(\zeta,w)-m\log|\zeta|^2$ is subharmonic on the punctured disk and, by the pole condition, bounded above near zero, so it extends as a subharmonic function across zero.

\Needspace{7\baselineskip}
Since $u\leq0$, this extension is at most $-m\log r^2$ on every circle $|\zeta|=r<R$. The maximum principle and then $r\uparrow R$ give
\[
 u(\zeta,w)\leq m\log\frac{|\zeta|^2}{R^2}.
\]
On $D$ both sides are $-\infty$. Taking the supremum over $u$ proves the formula. Near a point of $D$, the function $H_m$ is bounded and the logarithmic term tends to $-\infty$, which proves continuity there as an extended-real function.
\end{proof}

For a finite continuous obstacle $\Phi$ on $B$, put
\[
 P_B\Phi=\left(\sup\{v\in\PSH(B):v\leq\Phi\}\right)^*.
\]
The local upper-envelope theorem used in \cref{sec:obstacle-envelopes} applies whenever this family is nonempty and locally uniformly bounded above. Continuity of $\Phi$ ensures that the regularized envelope still satisfies $P_B\Phi\leq\Phi$.

\Needspace{10\baselineskip}
\begin{proposition}\label{prop:ball-green-truncation-rate}\index{Envelope!logarithmic obstacle}
For $t>0$, define
\[
 \Phi_t(\zeta,w)=\min\{-m\log|\zeta|^2,t\},
 \qquad V_t=P_B\Phi_t,
\]
with $\Phi_t=t$ on $D$. Then $0\leq V_t\leq t$, and
\begin{equation}\label{eq:ball-green-truncation-squeeze}
 -m\log\bigl(1-(1-\mathrm e^{-t/m})|w|^2\bigr)
 \leq V_t(\zeta,w)\leq H_m(w).
\end{equation}
In particular, $V_t\uparrow H_m$ pointwise on all of $B$, including $D$. The convergence is uniform on every compact subset, with
\begin{equation}\label{eq:ball-green-truncation-rate}
 0\leq H_m-V_t
 \leq \frac{m\mathrm e^{-t/m}R^2}{1-R^2}
 \qquad\text{where }|w|\leq R<1.
\end{equation}
\end{proposition}
\begin{proof}
The obstacle is continuous, nonnegative and at most $t$. Zero is a competitor, so the upper-envelope theorem gives $V_t\in\PSH(B)$ with $0\leq V_t\leq\Phi_t\leq t$. Consequently $m\log|\zeta|^2+V_t$ is a competitor for $G_m$, and \cref{prop:ball-divisor-green} gives $V_t\leq H_m$ off $D$. At a point $(0,w)$, the submean inequality on any sufficiently small circle in the normal line gives the same bound, since $H_m$ does not depend on $\zeta$.

For the lower bound, set $a_t=1-\mathrm e^{-t/m}$ and $F_t(w)=-m\log(1-a_t|w|^2)$. The Hessian calculation in the proof of \cref{prop:ball-divisor-green}, with $|w|^2$ replaced by $a_t|w|^2$, shows that $F_t$ is plurisubharmonic. Moreover, $0\leq F_t\leq t$, and $F_t\leq H_m\leq-m\log|\zeta|^2$ on $B\setminus D$ because $|\zeta|^2<1-|w|^2$. Thus $F_t\leq\Phi_t$ everywhere, and hence $F_t\leq V_t$. Finally,
\[
 H_m-F_t
 =m\log\left(1+\frac{\mathrm e^{-t/m}|w|^2}{1-|w|^2}\right)
 \leq\frac{m\mathrm e^{-t/m}|w|^2}{1-|w|^2}.
\]
This proves \eqref{eq:ball-green-truncation-rate} and the convergence; the monotonicity of $V_t$ in $t$ follows from that of $\Phi_t$.
\end{proof}

\Needspace{9\baselineskip}
\begin{example}\label{ex:ball-green-boundary-clusters}\index{Green function!boundary cluster set}
At every point $p=(0,w_0)\in\partial B$ with $|w_0|=1$, the cluster set of $G_m$ at $p$, that is, the set of limits of $G_m$ along sequences in $B$ converging to $p$, is all of $[-\infty,0]$. Indeed, set $w_s=(1-s^2)^{1/2}w_0$ for $0<s<1$. For every fixed $0<a<1$, the points $(as,w_s)$ belong to $B$, tend to $p$ as $s\to0$, and satisfy $G_m=m\log a^2$. Choosing $a=a_s<1$ with $a_s\to1$ gives the limit zero, whereas the points $(s^2,w_s)$ give $G_m=m\log s^2\to-\infty$. These realize every value in $[-\infty,0]$, and negativity excludes all others. Thus $G_m$, although continuous on $B$, has no continuous extension to any point of $\partial B$ in the closure of $D$.
\end{example}
